\par\Needspace{6\baselineskip}\medskip
\hypertarget{native-099-article076}{}\pdfbookmark[2]{76.}{article-076}
\begin{center}76.\end{center}
In art.\ 52 variationes differentiales longitudinis et latitudinis heliocentricae distantiaeque curtatae per variationes argumenti latitudinis $u$, inclinationis $i$ radiique vectoris $r$ exprimere docuimus, posteaque (art.\ 64, IV) ex illis deduximus variationes longitudinis et latitudinis geocentricae, $l$ et $b$: per combinationem itaque harum formularum $\mathrm{d}l$ et $\mathrm{d}b$ per $\mathrm{d}u$, $\mathrm{d}i$, $\mathrm{d}\AscendingNode$, $\mathrm{d}r$ expressae habebuntur. Sed operae pretium erit ostendere, quomodo in hoc quoque calculo reductione loci heliocentrici ad eclipticam supersedere liceat, sicuti in art.\ 65 locum geocentricum immediate e loco heliocentrico in orbita deduximus. Ut formulae eo simpliciores evadant, latitudinem terrae negligemus, quum certe in formulis differentialibus effectum sensibilem habere nequeat. Praesto sunt itaque formulae sequentes, in quibus brevitatis caussa $\omega$ pro $l-\AscendingNode$, nec non ut supra $\Delta'$ pro $\Delta\cos b$ scribimus:
\[
\begin{aligned}
\Delta'\cos\omega &= r\cos u - R\cos(L-\AscendingNode) &&= \xi\\
\Delta'\sin\omega &= r\cos i\sin u - R\sin(L-\AscendingNode) &&= \eta\\
\Delta'\operatorname{tang}b &= r\sin i\sin u &&= \zeta,
\end{aligned}
\]
\noindent e quarum differentiatione prodit
\SourceLeaf{100}\[
\begin{aligned}
\cos\omega\mathbin{.}\mathrm{d}\Delta' - \Delta'\sin\omega\mathbin{.}\mathrm{d}\omega &= \mathrm{d}\xi\\
\sin\omega\mathbin{.}\mathrm{d}\Delta' + \Delta'\cos\omega\mathbin{.}\mathrm{d}\omega &= \mathrm{d}\eta\\
\operatorname{tang}b\mathbin{.}\mathrm{d}\Delta' + \frac{\Delta}{\cos b}\mathrm{d}b &= \mathrm{d}\zeta.
\end{aligned}
\]
\noindent Hinc per eliminationem
\[
\begin{aligned}
\mathrm{d}\omega &= \frac{-\sin\omega\mathbin{.}\mathrm{d}\xi+\cos\omega\mathbin{.}\mathrm{d}\eta}{\Delta'}\\
\mathrm{d}b &= \frac{-\cos\omega\sin b\mathbin{.}\mathrm{d}\xi-\sin\omega\sin b\mathbin{.}\mathrm{d}\eta+\cos b\mathbin{.}\mathrm{d}\zeta}{\Delta}\mathpunct{\cdot}
\end{aligned}
\]

Si in his formulis pro $\xi$, $\eta$, $\zeta$ valores sui rite substituuntur, $\mathrm{d}\omega$ et $\mathrm{d}b$ per $\mathrm{d}r$, $\mathrm{d}u$, $\mathrm{d}i$, $\mathrm{d}\AscendingNode$ expressae prodibunt; dein, propter $\mathrm{d}l=\mathrm{d}\omega+\mathrm{d}\AscendingNode$, differentialia partialia ipsarum $l$, $b$ ita se habebunt:
\[
\mathrm{I.}\quad \Delta'\left(\frac{\mathrm{d}l}{\mathrm{d}r}\right) = -\sin\omega\cos u + \cos\omega\sin u\cos i
\]
\[
\mathrm{II.}\quad \frac{\Delta'}{r}\left(\frac{\mathrm{d}l}{\mathrm{d}u}\right) = \sin\omega\sin u + \cos\omega\cos u\cos i
\]
\[
\mathrm{III.}\quad \frac{\Delta'}{r}\left(\frac{\mathrm{d}l}{\mathrm{d}i}\right) = -\cos\omega\sin u\sin i
\]
\[
\begin{aligned}
\mathrm{IV.}\quad \left(\frac{\mathrm{d}l}{\mathrm{d}\AscendingNode}\right) &= 1+\frac{R}{\Delta'}\cos(L-\AscendingNode-\omega)\\
& = 1+\frac{R}{\Delta'}\cos(L-l)
\end{aligned}
\]
\[
\begin{aligned}
\mathrm{V.}\quad \Delta\left(\frac{\mathrm{d}b}{\mathrm{d}r}\right) &= -\cos\omega\cos u\sin b\\
& - \sin\omega\sin u\cos i\sin b + \sin u\sin i\cos b
\end{aligned}
\]
\[
\begin{aligned}
\mathrm{VI.}\quad \frac{\Delta}{r}\left(\frac{\mathrm{d}b}{\mathrm{d}u}\right) &= \cos\omega\sin u\sin b\\
& - \sin\omega\cos u\cos i\sin b + \cos u\sin i\cos b
\end{aligned}
\]
\[
\mathrm{VII.}\quad \frac{\Delta}{r}\left(\frac{\mathrm{d}b}{\mathrm{d}i}\right) = \sin\omega\sin u\sin i\sin b + \sin u\cos i\cos b
\]
\[
\begin{aligned}
\mathrm{VIII.}\quad \frac{\Delta}{R}\left(\frac{\mathrm{d}b}{\mathrm{d}\AscendingNode}\right) &= \sin b\sin(L-\AscendingNode-\omega)\\
& = \sin b\sin(L-l).
\end{aligned}
\]
\noindent Formulae IV et VIII hic iam in forma ad calculum commodissima apparent; formulae I, III, V autem per substitutiones obvias ad formam concinniorem rediguntur, puta
\[
\mathrm{I}^{*}.\quad \left(\frac{\mathrm{d}l}{\mathrm{d}r}\right) = \frac{R}{r\Delta'}\sin(L-l)
\]
\[
\mathrm{III}^{*}.\quad \left(\frac{\mathrm{d}l}{\mathrm{d}i}\right) = -\cos\omega\operatorname{tang}b
\]
\[
\begin{aligned}
\mathrm{V}^{*}.\quad \left(\frac{\mathrm{d}b}{\mathrm{d}r}\right) &= -\frac{R}{r\Delta}\cos(L-l)\sin b\\
& = -\frac{R}{r\Delta'}\cos(L-l)\sin b\cos b.
\end{aligned}
\]

\SourceLeaf{101}Denique formulae reliquae quoque II, VI, VII per introductionem quorundam angulorum auxiliarium in formam simpliciorem abeunt: quod commodissime fit sequenti modo. Determinentur anguli auxiliares $M$, $N$ per formulas
\[
\begin{gathered}
\operatorname{tang}M = \frac{\operatorname{tang}\omega}{\cos i},\\
\operatorname{tang}N = \sin\omega\operatorname{tang}i = \operatorname{tang}M\cos\omega\sin i.
\end{gathered}
\]
\noindent Tunc simul fit
\[
\begin{aligned}
\frac{\cos M^2}{\cos N^2}&=\frac{1+\operatorname{tang}N^2}{1+\operatorname{tang}M^2}\\
&=\frac{\cos i^2+\sin\omega^2\sin i^2}{\cos i^2+\operatorname{tang}\omega^2}=\cos\omega^2:
\end{aligned}
\]
\noindent iam quum ambiguitatem in determinatione ipsorum $M$, $N$ per tangentes suas remanentem ad lubitum decidere liceat, hoc ita fieri posse patet, ut habeatur $\dfrac{\cos M}{\cos N}=+\cos\omega$, ac proin $\dfrac{\sin N}{\sin M}=+\sin i$. Quibus ita factis, formulae II, VI, VII transeunt in sequentes:
\[
\mathrm{II}^{*}.\quad \left(\frac{\mathrm{d}l}{\mathrm{d}u}\right) = \frac{r\sin\omega\cos(M-u)}{\Delta'\sin M}
\]
\[
\begin{aligned}
\mathrm{VI}^{*}.\quad \left(\frac{\mathrm{d}b}{\mathrm{d}u}\right)&=\frac{r}{\Delta}\left\{\cos\omega\sin i\cos(M-u)\cos(N-b)\right.\\
&\qquad\left.+\sin(M-u)\sin(N-b)\right\}
\end{aligned}
\]
\[
\mathrm{VII}^{*}.\quad \left(\frac{\mathrm{d}b}{\mathrm{d}i}\right) = \frac{r\sin u\cos i\cos(N-b)}{\Delta\cos N}\mathpunct{\cdot}
\]
\noindent Hae transformationes respectu formularum II, VII neminem morabuntur, respectu formulae VI autem aliqua explicatio haud superflua erit. Substituendo scilicet in formula VI primo $M-(M-u)$ pro $u$, prodit
\[
\begin{aligned}
\frac{\Delta}{r}\left(\frac{\mathrm{d}b}{\mathrm{d}u}\right) = {}&\cos(M-u)\left\{\cos\omega\sin M\sin b\right.\\
&\qquad\left.-\sin\omega\cos i\cos M\sin b+\sin i\cos M\cos b\right\}\\
&-\sin(M-u)\left\{\cos\omega\cos M\sin b\right.\\
&\qquad\left.+\sin\omega\cos i\sin M\sin b-\sin i\sin M\cos b\right\}.
\end{aligned}
\]
\noindent Iam fit $\cos\omega\sin M =\allowbreak \cos i^2\cos\omega\sin M +\allowbreak \sin i^2\cos\omega\sin M =\allowbreak \sin\omega\cos i\cos M +\allowbreak \sin i^2\cos\omega\sin M$; unde pars prior illius expressionis transit in
\[
\begin{aligned}
&\sin i\cos(M-u)\left\{\sin i\cos\omega\sin M\sin b + \cos M\cos b\right\}\\
&= \sin i\cos(M-u)\left\{\cos\omega\sin N\sin b + \cos\omega\cos N\cos b\right\}\\
&= \cos\omega\sin i\cos(M-u)\cos(N-b).
\end{aligned}
\]
\noindent Perinde fit $\cos N =\allowbreak \cos\omega^2\cos N +\allowbreak \sin\omega^2\cos N =\hfil\penalty-10000 \cos\omega\cos M +\allowbreak \sin\omega\cos i\sin M$; unde expressionis pars posterior transit in
\[
-\sin(M-u)\left\{\cos N\sin b - \sin N\cos b\right\} = \sin(M-u)\sin(N-b).
\]
\noindent Hinc expressio VI\textsuperscript{*} protinus demanat.

\SourceLeaf{102}Angulus auxiliaris $M$ etiam ad transformationem formulae I adhiberi potest, quo introducto assumit formam
\[
\mathrm{I}^{**}.\quad \left(\frac{\mathrm{d}l}{\mathrm{d}r}\right) = -\frac{\sin\omega\sin(M-u)}{\Delta'\sin M},
\]
\noindent e cuius comparatione cum formula I\textsuperscript{*} concluditur
\[
-R\sin(L-l)\sin M = r\sin\omega\sin(M-u);
\]
\noindent hinc etiam formulae II\textsuperscript{*} forma paullo adhuc simplicior tribui potest, puta
\[
\mathrm{II}^{**}.\quad \left(\frac{\mathrm{d}l}{\mathrm{d}u}\right) = -\frac{R}{\Delta'}\sin(L-l)\operatorname{cotang}(M-u).
\]

Ut formula VI\textsuperscript{*} adhuc magis contrahatur, angulum auxiliarem novum introducere oportet, quod duplici modo fieri potest, scilicet statuendo vel $\operatorname{tang}P=\dfrac{\operatorname{tang}(M-u)}{\cos\omega\sin i}$, vel $\operatorname{tang}Q=\dfrac{\operatorname{tang}(N-b)}{\cos\omega\sin i}$: quo facto emergit
\[
\begin{aligned}
\mathrm{VI}^{**}.\quad \left(\frac{\mathrm{d}b}{\mathrm{d}u}\right) &= \frac{r\sin(M-u)\cos(N-b-P)}{\Delta\sin P}\\
& = \frac{r\sin(N-b)\cos(M-u-Q)}{\Delta\sin Q}\mathpunct{\cdot}
\end{aligned}
\]
\noindent Ceterum quantitates auxiliares $M$, $N$, $P$, $Q$ non sunt mere fictitiae, facileque, quidnam in sphaera coelesti singulis respondeat, assignare liceret: quin adeo hoc modo aequationum praecedentium plures adhuc elegantius exhiberi possent per arcus angulosve in sphaera, quibus tamen eo minus hic immoramur, quum in calculo numerico ipso formulas supra traditas superfluas reddere non valeant.

